\section{Indication Summary}

\begin{tabularx}{\linewidth}{@{} c c c >{\raggedright\arraybackslash}p{22mm} >{\raggedright\arraybackslash}p{30mm} >{\raggedright\arraybackslash}X @{}}
  \toprule
  \rowcolor{primary!12}
  \textbf{ID} & \textbf{Method} & \textbf{Type} & \textbf{Size~(mm)} & \textbf{Location} & \textbf{Orientation / Remarks} \\
  \midrule
  IND-001 & MT & Linear & $3.2 \times 0.3$ & Zone~A, 12~o'clock,\newline 47~mm from flange face & Axial (0--5$^{\circ}$). Sharp, well-defined. Morphology consistent with grinding burn. \\
  IND-002 & MT & Rounded & $1.5$~dia. & Zone~A, 4~o'clock,\newline 112~mm from flange face & ---~~No linear tail. Likely subsurface inclusion exposed by acid strip. \\
  IND-003 & UT & Volumetric & $4.1 \times 2.8$ & Zone~C, upper fillet radius,\newline 9.0~mm depth & Radial. Signal 82\%~FSH at ref.\ + 6~dB. Probable forging inclusion cluster. \\
  IND-004 & UT & Planar & $6.5 \times 1.2$ & Zone~B, LH lug bore ID,\newline 3.2~mm depth & 30$^{\circ}$ to bore axis. Exceeds DAC by 4~dB. Consistent with fatigue initiation from pitting. \\
  \bottomrule
\end{tabularx}

\subsection{Indication Details}
\begin{enumerate}[leftmargin=*, itemsep=2pt, topsep=4pt]
  \item \textbf{IND-001 --- Linear MT Indication.} Orientation parallel to grinding lay
        (circumferential). Indication bleeds rapidly at 0.5~s dwell, indicating shallow depth
        ($<$~0.5~mm). Photographed at 10$\times$ magnification. Recommended for blend-out
        assessment per SRM~57-20-01.
  \item \textbf{IND-002 --- Rounded MT Indication.} Isolated, non-linear. Characteristic of
        non-metallic inclusion intersecting the surface after chemical stripping. Deemed
        non-relevant per BAC~5424 paragraph~7.2.3.
  \item \textbf{IND-003 --- Volumetric UT Indication.} Located at 9.0~mm below surface in the
        fillet radius transition zone. Signal morphology suggests a cluster of small reflectors
        rather than a single discrete flaw. Amplitude: 82\%~FSH (exceeds 50\%~DAC threshold).
  \item \textbf{IND-004 --- Planar UT Indication.} Oriented at approximately 30$^{\circ}$ to the
        bore centreline axis. Signal rises sharply and exceeds the DAC curve by 4~dB at the
        peak. Likely a fatigue crack initiating from bore-surface corrosion pitting.
\end{enumerate}

% ============ SECTION 4: ACCEPTANCE CRITERIA ============
